\documentclass{article}
\usepackage{amsmath,amssymb}
\usepackage{xurl}
\usepackage[hidelinks]{hyperref}
% Shared commands for notes in docs/paper-gaps/.

\DeclareUnicodeCharacter{2080}{\ifmmode{}_{0}\else\textsubscript{0}\fi}
\DeclareUnicodeCharacter{2081}{\ifmmode{}_{1}\else\textsubscript{1}\fi}
\DeclareUnicodeCharacter{2082}{\ifmmode{}_{2}\else\textsubscript{2}\fi}
\DeclareUnicodeCharacter{2099}{\ifmmode{}_{n}\else\textsubscript{n}\fi}

\newcommand{\Lean}{\textsc{Lean}}
\ExplSyntaxOn
\NewDocumentCommand{\leanid}{m}{
  \ifmmode
    \textsf{\detokenize{#1}}
  \else
    \group_begin:
    \urlstyle{sf}
    \tl_set:Nn \l_tmpa_tl {#1}
    \tl_replace_all:Nnn \l_tmpa_tl {\_} {_}
    \exp_args:NV \path \l_tmpa_tl
    \group_end:
  \fi
}
\ExplSyntaxOff
\providecommand{\tr}{\operatorname{tr}}
\newcommand{\tnleanIssueBase}{https://github.com/LionSR/TNLean/issues/}
\newcommand{\tnleanPrBase}{https://github.com/LionSR/TNLean/pull/}
\newcommand{\tnleanDocBase}{https://sirui-lu.com/TNLean/docs/}
\newcommand{\ghissue}[1]{\href{\tnleanIssueBase#1}{GitHub issue \##1}}
\newcommand{\ghpr}[1]{\href{\tnleanPrBase#1}{PR \##1}}
\newcommand{\doclink}[1]{\href{\tnleanDocBase#1.html}{\path{#1}}}

% Dirac notation and the identity operator, as in blueprint/src/macros/common.tex.
% \providecommand keeps note-local definitions authoritative.
\usepackage{dsfont}
\providecommand{\ket}[1]{|#1\rangle}
\providecommand{\bra}[1]{\langle #1|}
\providecommand{\braket}[2]{\langle #1 | #2 \rangle}
\providecommand{\ketbra}[2]{|#1\rangle\!\langle #2|}
\providecommand{\Id}{\mathds{1}}
\providecommand{\one}{\mathds{1}}
\providecommand{\C}{\mathbb{C}}
\providecommand{\MN}[1]{M_{#1}(\C)}
\providecommand{\MD}{M_D(\C)}

% Verdict marker: \gapnote{<kind>}{<status>}. Kind is the policy
% classification (clarification, local-correction, scope-restriction,
% unfaithful, false-source, open-gap); status is open, wip (elimination
% underway), resolved, or historical. Machine-read by the site index;
% prints a verdict line.
\newcommand{\gapnote}[2]{\par\noindent\textbf{Verdict:} #1 (#2).\par}

\title{Translation-Invariant MPS and the Trivial Phase under Circuits with
    Measurements: One Direction}
\author{}
\date{2026-10-03}
\begin{document}
\maketitle
\gapnote{scope-restriction}{open}

\section{What the source states}
The paragraph ``Phases of matter'' of~\cite{Piroli2021LOCC} compares sequences
of states $\Psi=\{\ket{\psi_M}\}_{M\ge M_0}$ and
$\Phi=\{\ket{\varphi_M}\}_{M\ge M_0}$ on chains of increasing size. One writes
$\Psi\mapsto\Phi$ if there are $k\in\mathbb N$ and states $\sigma_M$ such that
$\ket{\psi_M}$ is mapped to $\sigma_M$ by a composition of $k$ channels, each a
circuit of depth $f(M)$ assisted by measurements and classical communication,
and $\lVert\sigma_M-\ket{\varphi_M}\bra{\varphi_M}\rVert_1\to0$, with $f$
polylogarithmic in $M$. The sequences are equivalent if $\Psi\mapsto\Phi$ and
$\Phi\mapsto\Psi$. The theorem labelled \texttt{MPS\_classification} reads: ``In
1D, all translational invariant MPS with fixed bond dimension are in the same
phase as the trivial state.''

\section{What is formalized}
The formal statement proves the direction from the trivial sequence to the MPS
sequence.\footnote{Formal declaration:
    \leanid{MPSPreparation.isAsymptoticallyPreparedWithMeasurementsInDepth_normalizedMPVState}
    in \doclink{TNLean/MPS/Preparation/QCccClassification}.}
Unit vectors $\ket{\psi_N}$ are prepared deterministically from a product state
by a circuit with measurements followed by a circuit, of total depth at most
$C\log N$, and
$\lVert\ket{\psi_N}\bra{\psi_N}-\ket{\varphi_N}\bra{\varphi_N}\rVert_1\to0$,
where $\ket{\varphi_N}$ is the normalized periodic state. This is
$\Psi\mapsto\Phi$ with $k=2$ and $\sigma_M$ pure. The restrictions are these.
\begin{enumerate}
    \item \emph{One direction.} The converse direction, from the MPS sequence
          to the trivial one, needs channels acting on arbitrary input states,
          for example exchanging every site with a fresh ancilla in $\ket0$ and
          discarding the ancillas. The measurement model of the formalization
          prepares states from product states only, and this direction is not
          stated.
    \item \emph{Canonical form.} The tensor is taken in the canonical form
          $A^i=\bigoplus_j\operatorname{diag}(\mu_{j,1},\dots,\mu_{j,m_j})\otimes
              A^i_j$ of~\cite{Malz2024Preparation}, with every block normal in
          the gauge $\sum_iA_j^{i\dagger}A_j^i=1$, with a positive definite
          fixed point of trace one, and with the mixed transfer maps of
          distinct blocks of spectral radius below one. The source assumes
          the canonical form of normal tensors without loss of generality and
          writes its proof for one normal block; the formal statement does not
          derive the canonical form of an arbitrary tensor, nor the spectral
          condition on mixed transfer maps of blocks that are not related by a
          gauge and a phase. The same reading is recorded
          in~\cite{gap:mswc24_block_form_mixed_overlap}.
    \item \emph{Nonvanishing states.} The periodic states are assumed nonzero
          for all large $N$, as a sequence of states presupposes; compare
          \cite{gap:mswc24_depth_upper_bound_nonzero_state}.
\end{enumerate}

\section{Elimination plan}
Extend the measurement model to channels acting on an arbitrary input state of
the chain together with ancillas that are traced out, prove the direction from
the MPS sequence to the trivial one, and derive the canonical-form hypotheses
from the canonical form of an arbitrary translation-invariant tensor.

\bibliographystyle{alpha}
\bibliography{references}
\end{document}
